\documentclass[sigconf,anonymous,balance=true]{acmart}
\setcopyright{none}
\settopmatter{printacmref=true,printfolios=true}
\acmConference[ACE '27]{Australasian Computing Education Conference}{February 1--5, 2027}{Canberra, Australia}
\acmYear{2027}\copyrightyear{2027}
\acmDOI{}\acmISBN{}\acmPrice{}
\hypersetup{pdfauthor={},pdfcreator={LaTeX},pdfproducer={LaTeX}}
\usepackage{booktabs}
\emergencystretch=1em
\AtBeginEnvironment{thebibliography}{\interlinepenalty=10000}
\begin{document}

\title[From Research-Paper Links to Defensible Claims]{From Research-Paper Links to Defensible Claims: A PaperWeave Teaching Resource}
\author{Anonymous Author(s)}
\begin{abstract}
A research-paper link can denote component reuse, a conditional experimental choice, or an analytical comparison. Treating these meanings as interchangeable can turn a reasonable connection into an unsupported claim. We present PaperWeave as a computing-education resource for teaching students to inspect and repair such claims. A Year 11 classroom account involves 22 students across two sessions and 20 usefulness responses (mean 8.1/10, median 8). A reported difficulty concerns interpreting dataset-dependent retrieval as universal. Building on that account, we completed a source-grounded audit of all nine prototype links and developed 18 contrasting claim cards, source locators, answer rationales and a teacher feedback guide. The resource distinguishes supported, refuted and not-established claims; its judgments retain component roles, experimental conditions and the difference between comparison and inheritance. Worked feedback examples distinguish a correct label from a justified explanation. Four additional tasks use a paper outside the original map to make a prospective check of source-scoped reasoning concrete. Versioned locators make the teaching judgments inspectable; learning effectiveness remains unmeasured. The contribution is an adoptable, computing-specific explanation activity with a complete worked corpus and bounded practitioner reflection. The new materials were prepared after the classroom sessions; no comparative or lasting learning benefit is claimed.
\end{abstract}
\ccsdesc[500]{Social and professional topics~Computing education}
\keywords{computing education, AI literacy, research literacy, evidence evaluation, practitioner experience}
\maketitle

\section{Educational problem and contributions}
Research-based computing lessons often use diagrams to connect architectures, datasets and methods. Such representations make a field easier to navigate, but the meaning of an arrow is not self-evident. A learner may infer that a component is always required when a paper reports using it in one experiment, or infer historical influence from a conceptual similarity. Reading the relationship label and evaluating its warrant are different activities.

We address this problem through PaperWeave, a small browser-based map of six natural-language-processing (NLP) papers. The teaching objective is to explain a relationship using an appropriate source and retain conditions that limit the claim. The learner is not expected to reproduce the models or master their mathematical details. Instead, the activity uses concrete computing examples to practise a distinction familiar in scientific reasoning: what was actually implemented, what could be implemented, and what the evidence does not establish.

This practitioner account draws on two classroom sessions, descriptive feedback, a subsequent structural audit, and inspection of the underlying research papers. These sources answer different questions. Classroom feedback describes participants' perceptions; the audit describes the current artifact; source inspection motivates teaching examples. We do not combine them into a claim that PaperWeave caused an improvement in understanding.

The paper contributes three connected resources. \textbf{C1: a structured claim record} separates a link's meaning from the evidence for a particular claim, retaining its component role, scope and source context. \textbf{C2: a complete worked corpus} covers all nine prototype links with 18 contrasting claims, source locators, judgments and repair rationales; its traceability is checked against version-identified PDFs. \textbf{C3: an adoption and feedback package} translates the classroom difficulty and source analysis into a student response sheet, a staged lesson, worked feedback examples and four new-source tasks. These make the required reasoning visible beyond a final judgment label. The contribution lies in the completed computing-specific materials and their connection to practice. Claim verification, scientific explanation and graph navigation are established approaches.

\section{Relationship to prior work}
\subsection{A bounded objective within AI literacy}
Long and Magerko~\cite{long} identify competencies for understanding and critically evaluating AI. Touretzky et al.~\cite{touretzky} articulate a K--12 agenda for communicating AI concepts, limitations and implications. Ng et al.~\cite{ng} distinguish understanding, application, evaluation and ethical dimensions, while Mannila et al.~\cite{mannila} examine K--12 AI literacy through stakeholder perspectives at ACE. PaperWeave addresses a narrower target within this landscape: explaining what the evidence supports about a relationship between computing methods. Recognising model names or operating a graph is not sufficient evidence of that capability.

Age and prior experience also affect the design. Druga et al.~\cite{druga} studied younger children's encounters with AI across different backgrounds, highlighting variation in experience and collaboration. Their population differs from our Year 11 cohort; the relevant design lesson is to establish prerequisites rather than assume a uniform starting point. Our materials consequently separate technical vocabulary from judgments about a claim's scope.

\subsection{Source evaluation has a substantial educational history}
Wineburg~\cite{wineburg} studied differences between historians and school students in evaluating documentary evidence. Britt and Aglinskas~\cite{britt} developed and evaluated instruction in sourcing, contextualisation and corroboration. Goldman et al.~\cite{goldman} connect evaluation with navigation, monitoring and sense-making when learning from multiple Internet sources. These studies establish that consulting documents and reasoning with evidence are distinct skills requiring instructional support.

McGrew et al.~\cite{mcgrew} assessed students' evaluation of online claims, sources and evidence. PaperWeave is not a general web-credibility intervention: students work with curated scholarly papers and examine what a particular passage licenses them to say. The extension to computing is concrete. A cited method may be one experimental choice rather than an architectural requirement; a dataset qualifier can change a universal claim; a conceptual comparison does not establish historical influence. We contribute worked tasks for these distinctions, without claiming to originate source evaluation or to demonstrate transfer from other disciplines.

\subsection{Citation intent, claim verification and explanation}
Citation intent and claim support answer different questions. SciCite~\cite{scicite} classifies the purpose of a citation, including method use and comparison; that purpose alone does not establish the scope of a learner's paraphrase. SciFact~\cite{scifact} pairs scientific claims with evidence, support/refutation judgments and rationales, explicitly avoiding a global truth verdict from a bounded corpus. We adapt this distinction to teaching: a link may be legitimate while a stronger statement about it is refuted or not established. Our authored cards are instructional examples, not a new expert-annotated verification benchmark.

McNeill et al.~\cite{mcneill} make claim, evidence and reasoning explicit in school science and examine performance when explanation scaffolds are removed. Their chemistry setting differs from ours, and their findings do not validate the present activity. We use the precedent to require an explanation of why a passage warrants a claim and to distinguish supported classroom performance from later independent reasoning. The addition here is a worked computing corpus in which parameter sharing, update policy, dataset conditions and architectural necessity can change the judgment.

\subsection{Scholarly-reading interfaces are close precedents}
PaperWeaver~\cite{paperweaver} contextualises recommended papers in relation to a researcher's collection using citation-based descriptions and generated comparisons. Despite the similar name, PaperWeave is a separate classroom prototype. Synergi~\cite{synergi} supports scholarly synthesis through citation contexts and mixed-initiative interaction. These precedents make an inter-paper map or explanation of a connection insufficient as a novelty claim.

Other interfaces support important parts of the reading process. CiteRead~\cite{citeread} integrates localised commentary from citing papers; CiteSee~\cite{citesee} uses reading history to contextualise citations and aid discovery. Scim~\cite{scim} assists skimming, and ScholarPhi~\cite{scholarphi} provides definitions of terms and symbols where readers need them. These functions inform the proposed source-location and glossary support. PaperWeaver also analyses facts and relationships in participants' notes~\cite{paperweaver}; it would therefore be inaccurate to claim that prior interfaces only support navigation. Our distinction is the deliberately contrasting claim task and its secondary-computing adoption materials, not the invention of evaluating written understanding.

The educational difference is therefore the required learner product: a claim with its component role, supporting location, limiting conditions and a justified decision to retain, revise or withhold it. The same product can be elicited from a printed map. PaperWeave supplies a navigable setting for the activity; its educational contribution does not rely on inventing graph navigation, citation context or automated scholarly synthesis. Table~\ref{tab:position} summarises these task distinctions.

\begin{table*}[t]
\caption{Positioning by the task and evidence reported in prior work. This is not a performance comparison or an exhaustive feature inventory.}
\label{tab:position}
\begin{tabular}{@{}p{0.22\textwidth}p{0.34\textwidth}p{0.38\textwidth}@{}}
\toprule
Precedent & Task or output & How the present resource builds on it\\
\midrule
PaperWeaver; Synergi~\cite{paperweaver,synergi} & Contextualised paper comparisons and scholarly synthesis & Supply contrasting learner claims and repair rationales for a fixed secondary-computing activity.\\
SciCite~\cite{scicite} & Citation-purpose classification & Separate a relation's purpose from whether a particular stronger claim follows.\\
SciFact~\cite{scifact} & Claim--evidence judgments with rationales & Adapt source-relative judgment to component roles and conditions; no new verification algorithm is claimed.\\
Explanation scaffolds~\cite{mcneill} & Claims, evidence and reasoning; assessment with reduced support & Supply computing-specific prompts, source locations and feedback actions; later independent performance remains unmeasured.\\
\bottomrule
\end{tabular}
\end{table*}

\subsection{Representations, explanations and discussion}
Ainsworth's DeFT framework~\cite{ainsworth} relates the design and functions of multiple representations to the tasks learners must perform. It motivates asking students to translate an edge into a qualified sentence instead of assuming that switching graph views produces understanding. Mayer and Moreno~\cite{mayer} analyse how instructional design can reduce unnecessary cognitive load; here, a small corpus and short contextualised passages are practical design choices, not evidence of a measured load reduction.

The ICAP framework~\cite{icap} provides a rationale for explanation and substantive interaction. Computing peer-instruction research~\cite{peerinstruction} separates discussion, teacher intervention and later individual performance, while Qiu et al.~\cite{qiu} offer an ACE precedent for structured evaluation of AI-assisted computing material. Our reported paired activity is teacher-facilitated practice, not a validated implementation of all these frameworks. They guide the proposed response sheet and later evaluation without supplying missing outcome evidence for this cohort.

\section{Prototype and teaching context}
\subsection{The artifact}
PaperWeave provides a curated map of BERT~\cite{bert}, T5~\cite{t5}, BART~\cite{bart}, Dense Passage Retrieval (DPR)~\cite{dpr}, Retrieval-Augmented Generation (RAG)~\cite{rag}, and Fusion-in-Decoder (FiD)~\cite{fid}. Nine manually specified links connect these papers. Learners can select architecture, learning-strategy and conceptual lenses, inspect relationship descriptions, and switch to a text-list view. The lenses expose different interpretations of the same collection rather than different experimentally assigned conditions.

Each link has a curator-assigned status and an informal source description. A displayed status of \emph{verified} is a prototype category, not a statement that an independent reviewer has validated the relationship. The prototype does not automatically ingest arbitrary papers or extract their relations. It also does not provide a built-in peer discussion, disagreement, or claim-repair interface. Paired discussion in this account was a classroom practice around the tool.

Keeping the corpus small made it possible to describe particular relationships and inspect their sources. It also constrains the contribution: the activity is an NLP teaching case, not a demonstration that the approach scales to any computing topic. A text-list view provides an alternative way to inspect the same curated relations, but no accessibility evaluation was conducted.

\subsection{Classroom setting and sequence}
The activity took place on 4 and 11 June 2026 with 22 Year 11 students aged 16--17 and one computing teacher. Approximately 40 minutes per session were reported for the classroom programme. The computing/AI-literacy topic concerned NLP, transformer architectures, retrieval-augmented generation, and relationships between research papers.

Students had encountered introductory computing and AI concepts but varied in their familiarity with academic research papers. The reported activity sequence introduced research relationships before emphasizing source inspection and claim revision. Teacher-facilitated paired discussion occurred, including during the first-session intervention described below. Both sessions concerned the same curated paper collection.

Students were allocated by chance to two groups of 11. Group A completed a structured predict--inspect--justify--revise activity. Group B used a static annotated relationship map with the same research papers and supporting materials, followed by conventional relationship-identification questions. Both groups participated in both sessions and received the teacher's explanation of the DPR--FiD example. Availability of an interactive prototype feature does not establish that both groups used it; in particular, the static-map activity should not be described as interactive exploration of all graph lenses.

This manuscript reports the activity design and practitioner experience rather than a treatment-effect estimate. The teacher explanation was common to both activities, and group-specific task outcomes are unavailable. The distinction between the two activities should therefore not be read as evidence that structured checking outperformed the static map.

The classroom records comprise relationship-inspection worksheets, task responses, questionnaire results, selected feedback and observation notes. This account uses practitioner-supplied aggregate questionnaire counts, selected comments and a teaching narrative. We do not report a complete worksheet analysis or infer misconception frequencies from selected examples.

\subsection{Feedback and reporting boundaries}
A questionnaire invited all 22 PaperWeave students and received 20 responses. The item asked: ``How useful was PaperWeave in helping you understand the relationships between computing research papers and evaluate the evidence supporting those relationships?'' The scale ran from 1 (not at all useful) to 10 (extremely useful). Its wording combines relationship understanding and evidence evaluation, so a result cannot distinguish separate perceived effects on those constructs.

The practitioner report supplies six selected feedback comments. Their wording has not been checked against original response records, so we paraphrase them rather than present verbatim participant quotations. We have not conducted a formal thematic analysis, quantified theme frequencies, or estimated coding agreement.

Ethics approval was obtained for the classroom research. Written parental or guardian consent and student assent were obtained before research data collection. Participation was voluntary, and students could take part in ordinary classroom activities without contributing research data. Records used participant identifiers, while personal information and consent documents were held separately with restricted access.

\section{Reported classroom experience}
\subsection{Reported perceived usefulness}
Table~\ref{tab:ratings} gives the supplied PaperWeave distribution. Its 20 ratings sum to 162, giving a mean of 8.1; both central observations are 8, giving a median of 8. Twenty of 22 invitees responded (90.9\%), and 15 respondents rated usefulness at least 8 (75\%). These calculations verify the aggregate arithmetic, not the underlying response collection.

\begin{table}[ht]
\centering
\caption{Practitioner-supplied PaperWeave usefulness ratings: 20 respondents from 22 invitees, scale 1--10.}
\label{tab:ratings}
\begin{tabular}{lrrrrrr}
\toprule
Rating & 1--5 & 6 & 7 & 8 & 9 & 10\\
\midrule
Count & 0 & 2 & 3 & 8 & 5 & 2\\
\bottomrule
\end{tabular}
\end{table}

The respondents generally judged the activity useful on the supplied item. The distribution pools feedback from the classroom activities; response counts and ratings are not available separately for the two groups. It therefore cannot establish the perceived benefit of the structured activity alone. Usefulness ratings also do not establish an improvement in reasoning, retention, independent reading, or confidence calibration. The account provides no matched baseline for this rating or delayed measure. Teacher explanations, original-source access and paired discussion are not isolated.

\begin{table*}[t]
\caption{Selected feedback, paraphrased from the classroom report. No frequencies or verbatim-quotation claims are implied.}
\label{tab:feedback}
\begin{tabular}{@{}p{0.29\textwidth}p{0.29\textwidth}p{0.36\textwidth}@{}}
\toprule
Feedback in the supplied account & Bounded interpretation & Suggested teaching response\\
\midrule
The map helped connect NLP models and papers. & Some learners valued a navigational overview. & Ask for a sentence explaining a selected link, to distinguish navigation from understanding.\\
Connected papers need not use the same method. & A learner noticed that links can have different meanings. & Contrast a component-use relation with a conceptual comparison using the same response form.\\
It was difficult to tell whether a relation was supported by its paper. & A short label did not settle the evidential question. & Model how to locate a passage and state what it supports, including conditions.\\
Different graph views supported different perspectives. & The lenses were perceived as useful for exploration. & Have learners explain why a relation appears under a particular lens.\\
Checking the original paper was useful. & Source inspection was positively described. & Require a source location alongside the student's claim.\\
NLP terminology required additional explanations. & Domain vocabulary could obstruct the intended reasoning task. & Supply a brief glossary and begin with one relation before opening the whole map.\\
\bottomrule
\end{tabular}
\end{table*}

\subsection{Selected feedback and interpretation}
The supplied examples describe both benefits and difficulties. They are selected feedback, rather than an exhaustive account of all responses. Table~\ref{tab:feedback} connects each paraphrase to a bounded interpretation and a possible teaching response. These responses are recommendations derived from the account; they are not claims that the teacher implemented each adaptation during the sessions.



The terminology difficulty is consequential for adoption. If students cannot distinguish a retriever from a generator, an incorrect relationship judgment may reflect unfamiliar vocabulary rather than weak source evaluation. We would therefore separate a short vocabulary check from the evidence task. Conversely, memorizing the names of the models should not earn credit for evaluating a claim without a source.

The feedback about support for relations highlights a productive tension. The map is useful as an overview, yet its concise labels can encourage reliance on the interface. Making that tension explicit turns the interface into an object of examination: learners should be able to explain why a displayed relation warrants qualification. However, the selected comments do not establish how many students made such a transition, or whether it persisted without the tool.

\subsection{A teacher-guided correction of an overbroad claim}
The practitioner describes a specific first-session difficulty on 4 June. Several students interpreted the DPR--FiD relationship as meaning that FiD always uses DPR. The teacher paused the activity and directed students to the original FiD paper, asking them to identify the retrieval methods used across datasets and decide whether one method was universal. Students inspected the evidence individually for several minutes, discussed interpretations with a partner, and rewrote the relationship to distinguish DPR for Natural Questions and TriviaQA from BM25 for SQuAD.

After the intervention, students returned to their worksheets and revised their interpretations. The follow-up asked them to identify the dataset-dependent relationship and explain why the earlier interpretation was too broad. The teacher supplied further explanation distinguishing the general architecture from the implementation evaluated in the paper. No numerical follow-up result is available for this account.

The narrative describes an opportunity for learners to inspect source passages and question a displayed relationship's scope. It does not establish a measured change in source-checking behaviour, independent transfer or retention. The teacher delivered the DPR--FiD explanation to the whole class, so both activity groups received that support. The episode therefore illustrates a shared teaching response rather than an outcome attributable solely to the structured activity.

This episode supports a concrete teaching reflection. An overview can help students begin exploring a technical field while encouraging them to read more into a short edge label than its evidence supports. The existing DPR--FiD record already includes a dataset-dependent qualifier. The account therefore does not demonstrate that the qualifier was absent; it suggests that learners may need guidance to notice and interpret it. A repeat lesson should make the conditional reading explicit, ask for a source-supported explanation, and distinguish the need for vocabulary support from a difficulty reasoning about scope.

\section{Source-grounded audit and completed task bank}
\subsection{What was examined and how}
The later artifact analysis covers all nine exported prototype links, rather than a sample chosen for favourable outcomes. The original records contain five \emph{verified}, three \emph{conceptual} and one \emph{inferred} status. All have informal source descriptions; none has explicit source-version, exact-passage or reviewer fields. These are representation properties, not estimates of relation correctness. In particular, the DPR--FiD source description already notes dataset dependence.

For each edge we inspected the cited primary paper in context, identified the strongest suitably scoped teaching claim, and recorded a contrasting claim that changes a consequential premise. We recorded document version, section, PDF page, a short locator anchor, the role or condition at issue, a judgment and a repair rationale. Two-document comparisons retain both sources. Fifteen locators span five source PDFs: DPR, BART, T5, RAG and FiD. BERT remains a map node, but claims about its use are checked in the paper that describes using or comparing it. The RAG locators refer to arXiv version 4; the other PDFs are the cited proceedings or journal versions. File hashes fix the inspected snapshots.

This source inspection and card authoring were AI-assisted. No independent human annotators established a gold standard. A separate script checks that every original edge has a pair of cards, every card's source identifiers resolve, and every locator anchor occurs on its declared PDF page. These checks test completeness and traceability; they cannot determine whether a rationale is semantically correct. The materials were completed after the June sessions and were not retrospectively administered or scored.

\subsection{Earlier review and resulting design distinctions}
The practitioner reports that an independent computing-education researcher reviewed an earlier prototype or draft and documented feedback on relationship accuracy, evidence, direction and scope. Three concerns were identified: distinguish DPR use by dataset in FiD; separate FiD's T5-based experiments from an architectural requirement; and avoid treating a T5--BART conceptual comparison as historical influence. These concerns are reflected in E6, E7 and E9 and make the design rationale more specific than a general request for clearer labels.

This feedback concerned an earlier version, not the completed 18-card bank or its answer key. The original feedback sheet was not available for inspection in preparing this account, and no review date or item-level agreement is reported. We therefore use it as practitioner-reported design feedback, not independent validation of the present judgments. A second content review of the current materials remains necessary.

\subsection{A record of the claim and its evidence}
A displayed edge is retained as an object of scrutiny. The teaching record contains: (1) the learner's asserted relationship; (2) the relevant component and relation meaning; (3) the condition or scope; (4) source version and location; and (5) a justified decision to retain, revise or withhold the claim. Relation meaning is distinct from judgment. A conceptual comparison can be supported, while a stronger statement about a component-use link can be refuted.

We use three source-relative judgments. \emph{Supported} means the identified context warrants the scoped claim. \emph{Refuted} requires an explicit incompatible statement or counterexample. \emph{Not established} means the identified evidence is insufficient; it does not declare the claim globally false. This follows established claim-verification reasoning~\cite{scifact}. The educational work is in identifying which premise changes the judgment and repairing the explanation, rather than merely selecting a label.

\subsection{Complete coverage and worked findings}
Table~\ref{tab:claims} provides one supported statement and one challenge for every edge. The nine pairs yield 18 authored cards: nine supported, seven refuted and two not established. These are design counts, not a measured classifier accuracy or distribution of student errors. The source analysis distinguishes three component-use relationships, one conditional use, one experimental instantiation, three author-stated conceptual relationships and one curator comparison.

\begin{table*}[t]
\caption{The complete nine-pair teaching bank. S is the supported claim; R and N denote refuted and not established challenges. Judgments are relative to the named source context. Wording is authored for teaching, not a transcript of interface labels or student answers.}
\label{tab:claims}
\begin{tabular}{@{}p{0.12\textwidth}p{0.29\textwidth}p{0.31\textwidth}p{0.22\textwidth}@{}}
\toprule
Pair & Supported claim (S) & Contrasting challenge and judgment & Source locator\\
\midrule
E1 BERT--DPR & DPR uses independent BERT encoders for questions and passages. & One parameter set is shared by both encoders. \textbf{R} & DPR \S3.1, p.\ 3~\cite{dpr}\\
E2 BERT--BART & BART compares its bidirectional encoder with BERT. & BART initializes by loading a BERT checkpoint. \textbf{R} & BART abstract; \S2.1, pp.\ 1--2~\cite{bart}\\
E3 BERT--T5 & The baseline objective draws on BERT and word dropout. & BERT is its sole stated inspiration. \textbf{R} & T5 \S3.1.4, pp.\ 12--13~\cite{t5}\\
E4 DPR--RAG & RAG uses DPR initialization and freezes its document encoder. & RAG fine-tunes both retrieval encoders. \textbf{R} & RAG \S\S2.2, 2.4, pp.\ 3--4~\cite{rag}\\
E5 BART--RAG & The BART-large generator is fine-tuned. & Its generator stays frozen during fine-tuning. \textbf{R} & RAG \S\S2.3--2.4, pp.\ 3--4~\cite{rag}\\
E6 DPR--FiD & DPR is used for NQ/TriviaQA; BM25 for SQuAD. & DPR is used for every reported dataset. \textbf{R} & FiD \S4, p.\ 4~\cite{fid}\\
E7 T5--FiD & Experiments initialize from T5; the method allows BART. & The architecture requires T5 and excludes BART. \textbf{R} & FiD \S\S3--4, pp.\ 3--4~\cite{fid}\\
E8 RAG--FiD & FiD compares passage processing with RAG. & The comparison establishes that FiD loads a trained RAG system. \textbf{N} & FiD \S\S2--4, pp.\ 2--4~\cite{fid}\\
E9 T5--BART & Both admit an encoder-decoder, denoising comparison. & Shared properties establish T5's historical influence on BART. \textbf{N} & T5 \S\S3.1.1, 3.1.4; BART \S2~\cite{t5,bart}\\
\bottomrule
\end{tabular}
\end{table*}

\textbf{A method-use label does not settle what was reused.} DPR's independent encoders distinguish a shared architecture from shared parameters~\cite{dpr}. BART's architectural comparison is accompanied by a separate random-initialisation description~\cite{bart}. RAG's retriever and generator have different update policies~\cite{rag}. E1, E2, E4 and E5 therefore ask learners to identify the component and operation, not infer every implementation detail from an arrow.

\textbf{Locators must preserve consequential context.} T5's baseline inspiration statement continues across pages 12--13 and names both BERT and word dropout~\cite{t5}. FiD's dataset mapping and T5 initialisation appear in its experimental details, beyond the method-level source locator~\cite{fid}. These cases justify supplying a context span rather than an isolated keyword match. The completed resource resolves all 15 anchors, but that result is a mechanical traceability check; it does not validate the interpretations independently.

\textbf{Comparison and uncertainty remain legitimate outcomes.} E8 and E9 retain the supported comparison while withholding the stronger inheritance or historical claim. They deliberately avoid treating ``not shown here'' as ``false everywhere.'' This is particularly important for the original inferred T5--BART link, whose source text already requests review. A curator can preserve a useful comparison without upgrading its evidential status.

\subsection{What this adds to the classroom account}
The observed DPR--FiD overgeneralisation motivates the distinction between a conditional claim and a universal one. The completed corpus makes related distinctions inspectable across every current edge, rather than leaving them as future interface suggestions. It also exposes different teaching responses: vocabulary support for a role error, a counterexample for a universal claim, or an additional source for an unsupported historical assertion. These are reasoned adoption implications. The audit does not show that students made every listed error or that the new cards improve their learning.

\section{Materials for another educator}
\subsection{Student materials and answer rationale}
The completed resource contains a shuffled 18-card student workbook, a separate teacher key, a source manifest and machine-readable claim records. The student version omits judgments and rationales. Each response records an initial interpretation, source location, consequential condition, final judgment and justified revision or evidence request. Table~\ref{tab:claims} supplies the corpus within the paper, so the teaching idea does not depend on access to an external prototype.

For E6, an adequate response identifies SQuAD as a counterexample to universal DPR use and restores the dataset-specific statement. For E9, an adequate response preserves the analytical similarity but requests evidence of historical influence; merely answering ``false'' loses the uncertainty distinction. These worked rationales are authored materials, not quotations from worksheets.

The teacher should inspect the complete cited paragraphs, explain the vocabulary needed for the assigned cards, and keep the key separate until learners commit an answer. Source access is required, but running the NLP models is not. A printed map and workbook can replace the browser interface.

\subsection{Preparing and running an adapted lesson}
An educator can reproduce the core reasoning task with the six-paper map, copies of the relevant papers, the prompts in Table~\ref{tab:claims} and a blank response sheet. The browser prototype is optional for this adaptation. Before teaching, check the cited passage and retain its surrounding context, identify the relation's direction and component role, and decide which terminology learners need. Review the source again if its version changes. No model training, API account or automatic relation extraction is required.

For a new implementation, a practical core uses E1, E6, E7 and E9: component identity, dataset scope, experimental choice versus necessity, and comparison versus history. The other five pairs are extensions for learners ready to inspect training and initialization details. In a two-session adaptation, model one source-based explanation first and then ask learners to resolve contrasting claims with less prompting. Within each task, learners first write their interpretation, inspect the source, explain its warrant and revise or withhold the claim. A partner can challenge an omitted condition before the teacher closes the discussion. This is an adaptation guide, not a minute-by-minute reconstruction of the two reported sessions.

A completed response to the FiD prompt could record: \emph{claim}---FiD always uses DPR; \emph{source}---FiD's experimental setup; \emph{condition}---retrieval differs by dataset; \emph{revision}---DPR is used for Natural Questions and TriviaQA, while BM25 is used for SQuAD; \emph{reason}---one counterexample defeats the universal claim. This worked response is derived from the paper and classroom narrative, not quoted from an individual worksheet. It shows what another teacher can request and assess without relying on a graph's status label.

An adequate response should identify the relationship's meaning, locate pertinent evidence and preserve the relevant condition. A source citation without an explanation is incomplete, as is a plausible explanation supported by the wrong section. Where evidence is insufficient, a justified decision to withhold the claim is acceptable. These are formative criteria for the reusable task, not a retrospectively applied or validated scoring instrument.

\subsection{Feedback that diagnoses the reasoning difficulty}
We provide four feedback criteria, adapted to the computing task rather than a validated numeric scale: correct relation meaning, relevant source location, preserved condition and a justified judgment. These distinguish different needs even when two learners choose the same final label. Table~\ref{tab:feedbackguide} gives teacher actions. No earlier worksheet was scored with this guide.

\begin{table}[t]
\caption{Prospective feedback guide; no scores or learning gains are inferred.}
\label{tab:feedbackguide}
\begin{tabular}{@{}p{0.43\columnwidth}p{0.51\columnwidth}@{}}
\toprule
Response difficulty & Teacher follow-up\\
\midrule
Names a model but confuses its role & Ask which component encodes, retrieves or generates.\\
Cites a paper without pertinent evidence & Locate the paragraph and explain its relevance.\\
Drops a dataset or implementation condition & Request a counterexample or qualified rewrite.\\
Treats absent support as refutation & Ask what evidence would settle the claim.\\
\bottomrule
\end{tabular}
\end{table}

For example, ``FiD uses multiple retrievers'' notices a distinction but leaves E6's decisive condition unspecified. A stronger explanation identifies the dataset mapping and why it defeats the universal claim. These are constructed response examples for teacher calibration, not additional participant data. Following the rationale of explanation-scaffolding work~\cite{mcneill}, a later task can remove the response headings while preserving source access. Whether learners then retain the reasoning process requires fresh evidence.

\subsection{Worked explanations and a new-source extension}
The assessment package makes the difference between a judgment and its warrant explicit. For E6, the constructed answer ``refuted'' supplies a label but no evidence. ``FiD uses more than one retriever'' notices variation but omits the condition. A sufficient explanation identifies FiD \S4, names SQuAD/BM25 as the counterexample, preserves the DPR mapping for Natural Questions and TriviaQA, and explains why one counterexample defeats ``every dataset.'' These levels guide feedback without assigning an unvalidated numerical score. For E9, a justified ``not established'' is preferable to an unsupported claim that historical influence is false.

The original pairs are useful practice but can cue one another's answers. We therefore supply four additional, unpaired statements about the Transformer paper~\cite{transformer}, which is outside the six-node map (Table~\ref{tab:newsource}). This extension asks learners to inspect configuration rows, preserve experimental scope and distinguish association from an isolated causal effect. It is a completed prospective activity, not a measured transfer result. It targets source-scoped reasoning within a new paper; it does not assess construction of a new inter-paper map.

\begin{table*}[t]
\caption{New-source extension: authored tasks and teacher rationale, not student outcomes. Locations use the pinned Transformer proceedings PDF.}
\label{tab:newsource}
\begin{tabular}{@{}p{0.06\textwidth}p{0.35\textwidth}p{0.39\textwidth}p{0.14\textwidth}@{}}
\toprule
Task & Claim & Judgment and required reasoning & Location\\
\midrule
T1 & Base and big configurations both use six layers per encoder/decoder stack. & Supported: locate $N=6$ in both named rows; do not generalise to every variation. & Table 3, p.\ 9\\
T2 & The big configuration uses eight attention heads. & Refuted: its $h$ value is 16; eight belongs to the base row. & Table 3, p.\ 9\\
T3 & The authors report nearly identical results for learned and sinusoidal positional representations in their comparison. & Supported: preserve the reported comparison's scope; near-equivalence is not a universal guarantee. & \S3.5, p.\ 6; Table 3\\
T4 & The base--big BLEU difference establishes that more heads alone caused the improvement. & Not established: dimensions, dropout and training steps also differ; request an otherwise controlled comparison. & \S6.2, pp.\ 8--9\\
\bottomrule
\end{tabular}
\end{table*}

A separate student sheet omits the key and asks for a judgment, location and explanation without the original step-by-step response headings. A brief glossary explains head, configuration and BLEU; source access remains available. The teacher key states what would warrant each judgment and identifies three source anchors, all resolved on their declared pages in the pinned PDF. This is an integrity result, not semantic or educational validation.

For an evaluation, reserve whole paper pairs from practice rather than splitting contrasting twins between practice and assessment. Check prior familiarity with the new source, document any vocabulary help and teacher prompting, and assess the evidence and explanation separately from label accuracy. An independent rater should calibrate on examples outside the evaluated responses before rating without condition information. The new tasks make such an evaluation feasible to inspect and adapt; their difficulty, scoring reliability and sensitivity to learning remain unmeasured.

\subsection{Pair discussion and reflection}
The new workbook asks learners to commit an initial interpretation before inspecting the assigned source context. An educator can then invite comparison with a partner's explanation. Each partner should identify the evidence supporting their claim, ask what would make the claim too broad, and record a revision if warranted. A pair may legitimately remain uncertain when the available evidence is insufficient. Agreement should not substitute for source support.

The classroom account attributes structured prediction, inspection, justification and revision to Group A and describes teacher-facilitated paired discussion. It does not establish a built-in disagreement interface or a common structured protocol for both groups. The new paper materials capture revisions explicitly; embedding that record in the software and evaluating its use remain future work.

The completed workbook and feedback guide operationalise these fields on paper. A built-in peer interface remains a future software direction, and no inter-rater reliability or outcome estimate is available for the new resource.

\section{Limitations and next steps}
This account concerns one small cohort, one computing topic, two sessions and one teacher. Both activity groups received the teacher explanation, so the reported episode does not isolate a benefit of the structured activity. The account depends on practitioner-selected observations and feedback, which may overrepresent salient or favourable experiences. The usefulness item combines two perceived benefits, pooled ratings obscure any group-specific pattern, and selected comments do not support frequency claims. The follow-up is described qualitatively, without learning-gain or between-group estimates.

The nine-link census is complete only for this small, curated prototype. Its 18 claims are deliberately constructed and potentially cue their intended contrast; they are not a representative corpus of misconceptions, an automatic-extraction benchmark or a validated test. The earlier researcher feedback does not independently validate the current answer key, source interpretations or new-source extension. The new manifest supplies explicit provenance for the teaching resource, but the original interface remains unchanged. Source-anchor checks demonstrate traceability, not semantic correctness. The new materials were not evaluated with students, and their suitability depends on prior knowledge and teacher scaffolding.

A fuller evaluation of the reported two-activity design would need group-specific outcomes, comparable assessment demands, documented allocation and teacher-intervention procedures, and a justified analysis. A subsequent primary outcome could assess reasoning about fresh claims, with raters unaware of condition. Transfer to new paper pairs would test a stronger claim than success on a rehearsed item. Confidence and completion time could be secondary outcomes, but neither would replace reasoning quality. These are evaluation requirements, not a report of completed comparative results.

Generative AI (OpenAI ChatGPT/Codex) assisted literature discovery, code development, source inspection, instructional card and rationale authoring, and manuscript drafting and revision. Classroom figures are practitioner-supplied; the technical audit examines a separate artifact snapshot. No human responses were generated for the reported classroom results.

\section{Conclusion}
PaperWeave turns research-paper links into claims that learners can explain, qualify or withhold. The classroom episode identifies a concrete scope error; the completed nine-link audit and 18-card resource extend it into inspectable computing examples. Versioned source locators, contrasting judgments, worked feedback and a new-source exercise make the activity adoptable without inventing a new mapping or verification algorithm. The contribution is a completed teaching resource informed by bounded classroom reflection. Independent validation of the answer key and evaluation of learning and transfer remain necessary for stronger empirical claims.
\bibliographystyle{ACM-Reference-Format}
\bibliography{references}
\end{document}
